\chapter{Modèles de crédit structurels}\label{ch:rc:structural-credit-models}

Le 4 mai 2005, le cours de l'action General Motors a bondi de 18~\% après que Kirk Kerkorian a proposé de racheter
une part importante du capital de l'entreprise~: sa société Tracinda a alors lancé une offre portant sur un maximum de 28 millions d'actions à 31
dollars chacune. Le lendemain, Standard \& Poor's a abaissé de deux
crans la note de la dette de la société, la ramenant à BB, sous le seuil de l'investment grade, et les spreads obligataires et ceux des credit default swaps se sont élargis.
Les opérations qui pariaient sur le maintien de l'alignement mutuel des cours de la dette et
de l'action de la société ont perdu sur les deux tableaux, deux jours consécutifs. Le modèle qui sous-tend bon nombre de ces opérations est le
sujet de ce chapitre~: une entreprise fait défaut lorsque la valeur de ses actifs tombe sous le niveau de sa
dette, de sorte que les fonds propres constituent une option sur l'entreprise et que les cours de ses actions et de son
crédit sont deux visions d'une même valeur d'actifs. Le \omterm{def:rc:reduced-form-credit:reduced}{modèle à forme réduite} du chapitre~13
prend l'intensité de défaut comme donnée~; le \omterm{def:rc:structural-credit-models:structural}{modèle structurel} l'explique et le relie au
marché des actions, ce qui fait son utilité et ce qui a rendu mai 2005 si coûteux.

\section{Les fonds propres comme option d'achat sur l'entreprise}

\begin{definition}[Modèle structurel]\label{def:rc:structural-credit-models:structural}
Un \emph{modèle structurel}\index{modèle structurel} de crédit modélise la valeur $V_t$ des actifs
d'une entreprise et fait du défaut un événement défini par $V$ et par les engagements de l'entreprise~:
les actifs passent sous la dette à son échéance ou touchent une \omterm{def:rc:structural-credit-models:fp}{barrière de défaut}. Le défaut est alors
prévisible à partir de la valeur des actifs, et les cours des fonds propres et de la dette en découlent.
\end{definition}

\begin{definition}[Modèle de Merton]\label{def:rc:structural-credit-models:merton}
Le \emph{modèle de Merton}\index{modèle de Merton} suppose que les actifs suivent un mouvement
brownien géométrique de volatilité $\sigma_V$ et une seule dette zéro-coupon de valeur faciale $D$ échéant à $T$. À $T$
les actionnaires reçoivent $\max(V_T-D,0)$ et les créanciers $\min(V_T,D)$. Les fonds propres sont une
option d'achat européenne sur les actifs, de prix d'exercice égal à la dette, évaluée par la
formule de Black--Scholes, et la dette vaut les actifs moins les fonds propres~:
\[
E_0 = V_0N(d_1)-De^{-rT}N(d_2),\qquad B_0 = V_0-E_0,\qquad
d_{1,2} = \frac{\ln(V_0/D)+(r\pm\tfrac12\sigma_V^2)T}{\sigma_V\sqrt T}.
\]
\end{definition}

De manière équivalente, la dette est une obligation sans risque diminuée d'une option de vente sur les actifs, de prix d'exercice $D$~: les
créanciers ont vendu aux actionnaires le droit de leur remettre l'entreprise au lieu de rembourser.
Le spread de crédit de la dette vaut $s = -\ln(B_0/D)/T - r$, et la \omterm{def:rc:reduced-form-credit:pd}{probabilité de défaut} risque-neutre
est $N(-d_2)$. La \cref{fig:rc:structural-credit-models:payoff} montre
les deux créances en fonction de la valeur des actifs.

\begin{omfigure}
\begin{tikzpicture}
\begin{axis}[width=11.5cm,height=5.4cm,
  xlabel={valeur des actifs $V_0$ (milliards de dollars)},ylabel={valeur (milliards de dollars)},
  tick label style={font=\small},label style={font=\small},
  xmin=0,xmax=80,ymin=0,ymax=45,grid=major,grid style={gray!25},
  legend columns=3,legend style={font=\footnotesize,at={(0.5,-0.25)},anchor=north,draw=none,fill=none,/tikz/every even column/.append style={column sep=8pt}},
  table/col sep=comma]
\addplot[omThm,very thick] table[x=v,y=equity]{figdata/rates-credit-risk/14-structural-credit-models/equity_debt.csv};
\addplot[omDef,very thick] table[x=v,y=debt]{figdata/rates-credit-risk/14-structural-credit-models/equity_debt.csv};
\addplot[omExo,thick,dashed] table[x=v,y=payoff]{figdata/rates-credit-risk/14-structural-credit-models/equity_debt.csv};
\legend{fonds propres aujourd'hui,dette aujourd'hui,fonds propres à l'échéance}
\end{axis}
\end{tikzpicture}

\omcaption{Fonds propres et dette de l'entreprise du chapitre (valeur faciale de la dette de 40~milliards de dollars échéant dans cinq
ans, volatilité des actifs de 15.6~\%) en fonction de la valeur des actifs. Les fonds propres sont une option d'achat sur les
actifs~; la dette croît vers la valeur sans risque de son remboursement et baisse d'un pour un avec les
actifs lorsqu'ils sont faibles. Données~: le tutoriel du chapitre.}\label{fig:rc:structural-credit-models:payoff}
\end{omfigure}

La valeur des actifs et leur volatilité ne sont pas observées. Les fonds propres le sont~: leur valeur de marché $E_0$
et leur volatilité $\sigma_E$, implicite des options ou estimée à partir des rendements. Le lemme d'Itô
appliqué à $E = C(V)$, $\sigma_EE_0 = N(d_1)\sigma_VV_0$, et avec l'équation de pricing, cela
donne deux équations pour les deux inconnues.

\omcode{code/firm/structural/firm_structural.py}{65}{81}{Retrouver la valeur et la volatilité des actifs à partir de la valeur et de la volatilité des fonds propres (Newton à deux inconnues).}\label{lst:rc:structural-credit-models:invert}

\begin{example}[Une entreprise endettée]\label{ex:rc:structural-credit-models:firm}
Une entreprise illustrative dispose de fonds propres valant 10~milliards de dollars, avec une volatilité de 50~\%, et une dette
de valeur faciale 40~milliards de dollars, considérée comme échéant dans cinq ans~; $r = 4~\%$. L'inversion donne des actifs
valant 40.69~milliards de dollars, avec une volatilité de 15.6~\%~: les fonds propres sont plus de trois fois
plus volatiles que les actifs, parce qu'elle est endettée. La dette vaut 30.69~milliards de dollars, soit un spread de rendement
de 130 points de base, et la probabilité risque-neutre que les actifs finissent sous la valeur faciale
est de 32.7~\%.
\end{example}

\section{Distance au défaut}

\begin{definition}[Distance au défaut]\label{def:rc:structural-credit-models:dd}
La \emph{distance au défaut}\index{distance au défaut} d'une entreprise sur un horizon $T$ est
le nombre d'écarts types dont la valeur espérée du logarithme des actifs dépasse le
point de défaut~:
\[
\mathrm{DD} = \frac{\ln(V_0/D)+(\mu-\tfrac12\sigma_V^2)T}{\sigma_V\sqrt T},
\]
avec $\mu$ la dérive historique des actifs. Dans le \omterm{def:rc:structural-credit-models:merton}{modèle de Merton}, la probabilité historique de
défaut est $N(-\mathrm{DD})$~; en pratique, la distance est convertie en fréquences de défaut
à partir d'un historique de défauts, car les rendements des actifs ne sont pas gaussiens.
\end{definition}

\begin{example}[Deux probabilités]\label{ex:rc:structural-credit-models:dd}
Avec une dérive des actifs de 8~\%, la \omterm{def:rc:structural-credit-models:dd}{distance au défaut} à cinq ans de l'entreprise vaut 1.02 et la
probabilité historique de défaut vaut 15.4~\%, contre les 32.7~\% risque-neutres de
l'\cref{ex:rc:structural-credit-models:firm}. L'écart est la \omterm{def:rc:reduced-form-credit:premium}{prime de risque de défaut}
du chapitre~13 sous forme structurelle~: sous la probabilité risque-neutre, les actifs dérivent au taux $r$
plutôt que $\mu$.
\end{example}

La structure par terme des spreads du \omterm{def:rc:structural-credit-models:merton}{modèle de Merton} dépend du quasi-levier
$d = De^{-rT}/V_0$ (la \cref{fig:rc:structural-credit-models:merton}). Une entreprise faiblement endettée
a des spreads qui partent de zéro et croissent avec l'échéance, car il faut du temps
pour que les actifs tombent jusqu'à la dette~; une entreprise proche du point de défaut a une courbe en cloche~; une
entreprise insolvable ($d>1$) a des spreads qui décroissent avec l'échéance, car seul le temps peut la sauver.

\begin{omfigure}
\begin{tikzpicture}
\begin{axis}[width=11.5cm,height=5.4cm,
  xlabel={échéance (années)},ylabel={spread de Merton (pb)},
  tick label style={font=\small},label style={font=\small},
  xmin=0,xmax=10,ymin=0,ymax=2000,grid=major,grid style={gray!25},
  restrict y to domain=0:2100,
  legend columns=3,legend style={font=\footnotesize,at={(0.5,-0.25)},anchor=north,draw=none,fill=none,/tikz/every even column/.append style={column sep=8pt}},
  table/col sep=comma]
\addplot[omMeth,very thick] table[x=t,y=d05]{figdata/rates-credit-risk/14-structural-credit-models/merton_ts.csv};
\addplot[omDef,very thick,dashed] table[x=t,y=d09]{figdata/rates-credit-risk/14-structural-credit-models/merton_ts.csv};
\addplot[omThm,very thick,dotted] table[x=t,y=d12]{figdata/rates-credit-risk/14-structural-credit-models/merton_ts.csv};
\legend{$d = 0.5$,$d = 0.9$,$d = 1.2$}
\end{axis}
\end{tikzpicture}

\omcaption{Spreads de rendement de Merton en fonction de l'échéance pour trois quasi-leviers
$d = De^{-rT}/V_0$, volatilité des actifs 25~\%~: croissant pour une entreprise solide, en bosse au voisinage du
point de défaut, décroissant pour une entreprise dont les actifs valent moins que sa dette actualisée. Données~: le
tutoriel du chapitre.}\label{fig:rc:structural-credit-models:merton}
\end{omfigure}

\begin{remark}[Spreads trop faibles]\label{rem:rc:structural-credit-models:low}
Les spreads de Merton sont faibles pour les entreprises solides et les échéances courtes~: une diffusion ne peut pas atteindre le
point de défaut en quelques mois, de sorte que les spreads à court terme sont proches de zéro alors que les marchés cotent des dizaines de
points de base. Pour y remédier, on ajoute des sauts aux actifs, une incertitude sur le point de défaut, ou
une barrière pouvant être atteinte à tout instant, objet de la section suivante.
\end{remark}

\section{Modèles de premier passage}

\begin{definition}[Modèle de premier passage, barrière de défaut]\label{def:rc:structural-credit-models:fp}
Un \emph{modèle de premier passage}\index{modèle de premier passage} fait défaut à l'entreprise la première fois que
ses actifs touchent une \emph{barrière de défaut}\index{barrière de défaut} $B_t$ située sous la
valeur courante, dès que cela se produit, et non seulement à l'échéance de la dette. Une barrière représente
des covenants permettant aux créanciers de prendre le contrôle, ou le niveau à partir duquel l'entreprise ne peut plus
se refinancer. Le modèle CreditGrades (2002), publié par RiskMetrics avec trois dealers, fixe la barrière
à la valeur moyenne du recouvrement sur la dette, multipliée par la dette par action, et rend ce
recouvrement incertain selon une loi lognormale, de sorte que le défaut peut survenir de façon inattendue et que les spreads à court terme ne sont pas nuls.
\end{definition}

\begin{proposition}[Survie à une barrière constante]\label{prop:rc:structural-credit-models:fp}
Avec une dynamique des actifs sous la probabilité risque-neutre $dV = rV\,dt+\sigma_VV\,dW$, une barrière constante $B<V_0$,
$x = \ln(V_0/B)$ et $\nu = r-\tfrac12\sigma_V^2$, la probabilité de non-défaut jusqu'à $t$ est
\[
Q(t) = N\!\left(\frac{x+\nu t}{\sigma_V\sqrt t}\right)
- e^{-2\nu x/\sigma_V^2}\,N\!\left(\frac{-x+\nu t}{\sigma_V\sqrt t}\right).
\]
\end{proposition}
\begin{proof}
$\ln V_t-\ln V_0$ est un mouvement brownien avec dérive $\nu$ et de volatilité $\sigma_V$~; la
probabilité qu'il reste au-dessus de $-x$ jusqu'à $t$ découle du principe de réflexion et
d'un changement de probabilité supprimant la dérive (livre~4, chapitre~2).
\end{proof}

\omcode{code/firm/structural/firm_structural.py}{93}{100}{La probabilité de survie de premier passage, sous forme fermée.}\label{lst:rc:structural-credit-models:fp}

La fonction de survie devient une courbe d'intensité de défaut aux dates de paiement de la
prime, et les jambes du chapitre~13 la transforment en spread de CDS du modèle pour chaque échéance, avec un recouvrement fixé selon la
convention du CDS (40~\%) plutôt que déduit de la valeur des actifs à la barrière.

\begin{example}[L'entreprise sous une barrière]\label{ex:rc:structural-credit-models:fp}
Avec la barrière à 70~\% de la valeur nominale de la dette, les spreads de CDS du modèle de l'entreprise sont de 64,
173, 216, 229 et 197 points de base à un, deux, trois, cinq et dix ans
(voir la \cref{fig:rc:structural-credit-models:fp})~; sa probabilité de survie à cinq ans est de 82.3~\%. Une
cotation de marché de 500 points de base à cinq ans exige une barrière à 78.7~\% de la valeur nominale.
\end{example}

\begin{omfigure}
\begin{tikzpicture}
\begin{axis}[width=11.5cm,height=5.4cm,
  xlabel={échéance (années)},ylabel={spread de CDS du modèle (pb)},
  tick label style={font=\small},label style={font=\small},
  xmin=0,xmax=10,ymin=0,ymax=650,grid=major,grid style={gray!25},
  legend columns=3,legend style={font=\footnotesize,at={(0.5,-0.25)},anchor=north,draw=none,fill=none,/tikz/every even column/.append style={column sep=8pt}},
  table/col sep=comma]
\addplot[omMeth,very thick,mark=*,mark size=1.2pt] table[x=t,y=l60]{figdata/rates-credit-risk/14-structural-credit-models/fp_ts.csv};
\addplot[omDef,very thick,mark=square*,mark size=1.2pt] table[x=t,y=l70]{figdata/rates-credit-risk/14-structural-credit-models/fp_ts.csv};
\addplot[omThm,very thick,mark=triangle*,mark size=1.6pt] table[x=t,y=limp]{figdata/rates-credit-risk/14-structural-credit-models/fp_ts.csv};
\legend{barrière à 60~\% de la valeur nominale,barrière à 70~\%,barrière à 78.7~\% (calée sur 500 pb)}
\end{axis}
\end{tikzpicture}

\omcaption{Spreads de CDS du modèle de l'entreprise du chapitre dans le \omterm{def:rc:structural-credit-models:fp}{modèle de premier passage},
pour trois barrières. Relever la barrière soulève la courbe et déplace sa bosse vers des
échéances plus courtes~; la barrière la plus élevée est calée sur une cotation à cinq ans de 500 points de base. Données~: le
tutoriel du chapitre.}\label{fig:rc:structural-credit-models:fp}
\end{omfigure}

\begin{dated}{2026-09}{Mai 2005}\label{dat:rc:structural-credit-models:may2005}
La revue de la Bank for International Settlements de juin 2005 rapporte que le 4 mai 2005
l'action de General Motors a progressé de 18~\% après une offre de Kirk Kerkorian pour acquérir une participation importante,
que le 5 mai Standard \& Poor's a abaissé la note de GM de deux crans, à BB, et celle de Ford d'un cran, à
BB+, et que les opérations de valeur relative entre les obligations convertibles de GM et un
portefeuille de réplication ont perdu sur les deux jambes, les spreads de l'obligation et du CDS ayant augmenté.
\end{dated}

\section{Signaux de structure du capital et arbitrage}

Un \omterm{def:rc:structural-credit-models:structural}{modèle structurel} relie un cours d'action à un spread de crédit. Les traders utilisent cette correspondance dans les deux
sens~: les mouvements de l'action prédisent ceux des spreads, et les spreads éloignés de la valeur du modèle
sont des opérations candidates.

\begin{definition}[Arbitrage de structure du capital]\label{def:rc:structural-credit-models:csa}
L'\emph{arbitrage de structure du capital}\index{arbitrage de structure du capital} oppose un
instrument de la structure du capital d'une entreprise contre un autre (CDS ou obligations
contre l'action ou les options sur action) lorsque leurs prix s'écartent de ceux que donne un \omterm{def:rc:structural-credit-models:structural}{modèle structurel},
en se couvrant grâce à la sensibilité du modèle de l'un à l'autre~: vendre la protection et vendre à découvert
l'action lorsque le spread de marché est supérieur à celui du modèle, acheter la protection et acheter
l'action lorsque celui-ci est inférieur.
\end{definition}
Yu (2006) a mis en œuvre une telle stratégie de convergence avec CreditGrades sur 135\,759 spreads quotidiens de CDS de 261
noms nord-américains~: des transactions isolées pouvaient perdre lourdement parce que les spreads et les cours des actions étaient faiblement
corrélés, alors qu'un portefeuille équipondéré de toutes les transactions affichait des ratios de Sharpe proches de ceux
d'autres stratégies d'arbitrage obligataire.


\begin{method}[La couverture par l'action d'un CDS]\label{met:rc:structural-credit-models:hedge}
Garder fixe la volatilité des actifs. Faire varier de peu la valeur des fonds propres, chercher la valeur
d'actif qui reproduit ce prix, et recalculer le spread du modèle~: cela donne $\partial s/\partial E$.
Multiplier par le \omterm{def:rc:reduced-form-credit:cs01}{CS01} de la position (par unité de spread) pour obtenir sa variation de valeur par unité de
valeur des fonds propres, et négocier autant d'actions dans le sens opposé.
\end{method}

\begin{example}[Le ratio de couverture]\label{ex:rc:structural-credit-models:hedge}
Avec la barrière à 70~\% de la valeur nominale, le spread du modèle baisse de 0.47 point de base pour
chaque hausse de 0.01~milliard de dollars de la valeur des fonds propres ($\partial s/\partial E = -0.0047$ par
milliard de dollars). 10~millions de dollars de protection vendue à cinq ans au coupon de 500 présente un \omterm{def:rc:reduced-form-credit:cs01}{CS01} de
$-\text{USD}~3\,713$ par point de base pour une cotation de 500, soit un gain de 174\,423~dollars par
milliard de dollars de valeur des fonds propres~: la couverture est une vente à découvert de 1.74~millions de dollars d'actions
(voir la \cref{fig:rc:structural-credit-models:equityspread}).
\end{example}

\begin{omfigure}
\begin{tikzpicture}
\begin{axis}[width=11.5cm,height=5.4cm,
  xlabel={valeur des fonds propres (milliards d'USD)},ylabel={spread du modèle à cinq ans (pb)},
  tick label style={font=\small},label style={font=\small},
  xmin=6,xmax=16,ymin=0,ymax=500,grid=major,grid style={gray!25},
  table/col sep=comma]
\addplot[omDef,very thick] table[x=e,y=spread]{figdata/rates-credit-risk/14-structural-credit-models/equity_spread.csv};
\addplot[omThm,only marks,mark=*,mark size=2.2pt] coordinates {(10,228.7) (11.8,160.7)};
 \node[omThm,font=\footnotesize,anchor=south west] at (axis cs:10,232) {avant~: 229~pb};
 \node[omThm,font=\footnotesize,anchor=south west] at (axis cs:11.8,164) {après +18~\%~: 161~pb};
\end{axis}
\end{tikzpicture}

\omcaption{Spread à cinq ans du \omterm{def:rc:structural-credit-models:fp}{modèle de premier passage} en fonction de la valeur des fonds propres de l'entreprise,
volatilité des actifs maintenue fixe (barrière à 70~\% de la valeur nominale). La pente à USD~10 milliards est le
ratio de couverture~; la convexité de la courbe explique pourquoi une variation de 18~\% laisse une position couverte en perte même
si le modèle est juste. Données~: le tutoriel du chapitre.}\label{fig:rc:structural-credit-models:equityspread}
\end{omfigure}

\begin{remark}[Ce que le modèle suppose sur la volatilité]\label{rem:rc:structural-credit-models:vol}
En maintenant la volatilité des actifs fixe, la volatilité des fonds propres baisse lorsque ceux-ci augmentent
(l'effet de levier diminue)~: après une hausse de 18~\%, elle vaut 47.1~\%. Si en revanche la volatilité des fonds propres reste à
50~\%, l'inversion donne une volatilité des actifs plus élevée, 17.3~\%, et le spread du modèle
\emph{augmente}, jusqu'à 251 points de base. Le signe de la couverture dépend de la volatilité que le
desk considère comme stable, et des options sur actions, si l'entreprise en émet.
\end{remark}

\section{Tutoriel~: des fonds propres aux spreads}\label{tut:rc:structural-credit-models:tutorial}

\begin{tutorial}
\textbf{Objectif.} Retrouver les actifs d'une entreprise à partir de ses fonds propres, mesurer la \omterm{def:rc:structural-credit-models:dd}{distance au défaut},
construire les spreads des \omterm{def:rc:structural-credit-models:merton}{modèles de Merton} et de premier passage, et dimensionner une couverture en actions.
\textbf{État final~:} les chiffres des
\cref{ex:rc:structural-credit-models:firm,ex:rc:structural-credit-models:fp,ex:rc:structural-credit-models:hedge}
et des \cref{fig:rc:structural-credit-models:payoff,fig:rc:structural-credit-models:merton,fig:rc:structural-credit-models:fp,fig:rc:structural-credit-models:equityspread}.
\begin{tutsteps}
\item \textbf{Inverser}~: \texttt{invert(10, 0.5, 40, 5, 0.04)}~; vérifier que les fonds propres et leur volatilité
sont bien recalculés.
\item \textbf{Merton}~: \texttt{firm()} pour le spread et les deux \omterm{def:rc:reduced-form-credit:pd}{probabilités de défaut}.
\item \textbf{Premier passage}~: \texttt{fp\_spread} pour chaque échéance et
\texttt{implied\_barrier()} pour la cotation à 500 points de base.
\item \textbf{Couverture}~: \texttt{hedge()}~; \texttt{fig\_rc\_structural.py} écrit les graphiques.
\end{tutsteps}
\textbf{Pour aller plus loin.} Faire croître la barrière au taux sans risque
(\texttt{growth}) et comparer les courbes~; recalculer la couverture avec la volatilité des fonds propres
maintenue constante au lieu de celle des actifs.
\end{tutorial}

\section{Construction~: modèles structurels}\label{bld:rc:structural-credit-models:structural}

\begin{build}
\bfield{Rôle} La cartographie des fonds propres vers le crédit d'une entreprise~: spreads de modèles pour les émetteurs à actions cotées,
la couverture en actions des positions de crédit, et un contrôle des valorisations de credit default swap (CDS).
\bfield{Interface} \texttt{merton\_equity}, \texttt{merton\_debt}, \texttt{merton\_spread},
\texttt{merton\_pd}, \texttt{distance\_to\_default}, \texttt{equity\_vol}~; \texttt{invert(E,
sigma\_E, D, T, r)}~; \texttt{asset\_from\_equity}~; \texttt{first\_passage\_survival}~;
\texttt{hazard\_curve\_from\_survival}~; \texttt{model\_cds\_spread(q, r, maturity,
recovery)}~; \texttt{FlatDisc}.
\bfield{Règles} Taux constant~; une seule valeur faciale de dette~; les fonctions de survie deviennent des courbes d'intensité de défaut sur la
grille trimestrielle et sont valorisées par \texttt{firm\_cdscurve} du chapitre~13.
\bfield{Tests d'acceptation} \texttt{code/firm/structural/tests/}~: l'inversion recalcule
les fonds propres et leur volatilité~; la dette plus les fonds propres égale les actifs, et la dette est inférieure à la
valeur sans risque~; la formule de premier passage se réduit au principe de réflexion sans dérive et à
celle d'une barrière nulle~; une barrière provoque un défaut plus souvent qu'un test terminal~; une
survie constante donne le spread au pair du chapitre~13.
\bfield{Pour aller plus loin} Barrière de niveau incertain (recouvrement aléatoire du point de défaut)~;
sauts dans les actifs~; une structure par terme de la dette~; la calibration, par les options sur actions, de la
volatilité des actifs.
\end{build}

\begin{omsources}
\item R.\ C.\ Merton, ``On the pricing of corporate debt: the risk structure of interest
rates'', \emph{Journal of Finance} 29(2), 1974, 449--470.
\item F.\ Black and J.\ C.\ Cox, ``Valuing corporate securities: some effects of bond
indenture provisions'', \emph{Journal of Finance} 31(2), 1976, 351--367.
\item F.\ Yu, ``How profitable is capital structure arbitrage?'', \emph{Financial Analysts
Journal} 62(5), 2006, 47--62.
\item C.\ C.\ Finger (éd.), \emph{CreditGrades Technical Document}, RiskMetrics Group, mai 2002.
\item Bank for International Settlements, \emph{Quarterly Review}, juin 2005, encadré ``Stress
testing of credit markets: the downgrade of General Motors and Ford''.
\end{omsources}

\section{Exercices}

\begin{exercise}[$\star$]\label{exo:rc:structural-credit-models:1}
Actifs de 50, dette de valeur faciale 40 à échéance de cinq ans, volatilité des actifs de 20~\%, $r = 4~\%$. Donner les
valeurs de Merton des fonds propres et de la dette, le spread de rendement et la \omterm{def:rc:reduced-form-credit:pd}{probabilité
de défaut} risque-neutre.
\end{exercise}

\begin{exercise}[$\star$]\label{exo:rc:structural-credit-models:2}
Pourquoi les spreads de Merton sont-ils proches de zéro aux courtes échéances pour une entreprise saine~?
\end{exercise}

\begin{exercise}[$\star$]\label{exo:rc:structural-credit-models:3}
Calculer la \omterm{def:rc:structural-credit-models:dd}{distance au défaut} à un an et la probabilité historique de défaut de
l'entreprise du chapitre, avec la même valeur faciale de dette et la même dérive.
\end{exercise}

\begin{exercise}[$\star\star$]\label{exo:rc:structural-credit-models:4}
Sans dérive, les actifs d'une entreprise valent 1.5 fois la barrière, avec une volatilité des actifs de 30~\%.
Calculer la probabilité qu'elle survive deux ans.
\end{exercise}

\begin{exercise}[$\star\star$]\label{exo:rc:structural-credit-models:5}
Dimensionner la couverture en actions de 20~millions d'USD de protection vendue sur l'entreprise du chapitre.
\end{exercise}

\begin{exercise}[$\star\star$]\label{exo:rc:structural-credit-models:6}
Pourquoi les fonds propres sont-ils beaucoup plus volatils que les actifs, et pourquoi le ratio change-t-il lorsque le
cours des fonds propres évolue~?
\end{exercise}

\begin{exercise}[$\star\star\star$]\label{exo:rc:structural-credit-models:7}
\emph{Programmation.} Trouver la barrière qui correspond à une cotation à cinq ans de 500 points de base, et donner
le spread du modèle à dix ans avec cette barrière.
\end{exercise}

\begin{exercise}[$\star\star\star$]\label{exo:rc:structural-credit-models:8}
\emph{Trouver la faille.} «~Le \omterm{def:rc:structural-credit-models:merton}{modèle de Merton} donne à cette entreprise 130 points de base et le marché
cote 500. Le credit default swap est trop large de 370 points de base~: il faut vendre la protection.~»
\end{exercise}

\section{Problème~: l'arbitrage de structure du capital}

\begin{problem}[{Problème du week-end --- mai 2005}]\label{pb:rc:structural-credit-models:1}
Un fonds constate que l'entreprise du chapitre est cotée à 500 points de base à cinq ans, alors que son
\omterm{def:rc:structural-credit-models:fp}{modèle de premier passage} (barrière à 70~\% du nominal, calibré sur l'historique de l'entreprise) donne 229. Il
vend pour 10~millions d'USD de protection à cinq ans au coupon de 500 et vend à découvert des actions selon le
ratio de couverture du modèle. Le lendemain, le cours de l'action augmente de 18~\%~; le surlendemain, l'entreprise est
déclassée, et le spread de marché s'élargit de 200 points de base, pour atteindre 700.

\medskip\noindent\textbf{Partie I --- L'opération.}
\begin{enumerate}
\item Donner le spread du modèle et le spread de marché, ainsi que le signal.
\item Donner le ratio de couverture et la vente à découvert sur l'action.
\item Quel gain le fonds réalise-t-il tant que rien ne se passe~?
\item Quel pari le fonds prend-il, énoncé sans recourir au modèle~?
\item Quels risques la couverture laisse-t-elle ouverts~?
\end{enumerate}

\medskip\noindent\textbf{Partie II --- Les deux jours.}
\begin{enumerate}[resume]
\item Donner le spread du modèle après la hausse de 18~\%.
\item Donner le spread de marché que prédit le modèle après la hausse, en appliquant
la variation du modèle à la cotation.
\item Donner le P\&L de la jambe actions.
\item Donner le P\&L du credit default swap si le marché avait suivi le modèle, ainsi que le P\&L couvert.
\item Donner le P\&L du credit default swap à 700 et le total.
\end{enumerate}

\medskip\noindent\textbf{Partie III --- Ce que le modèle n'a pas vu.}
\begin{enumerate}[resume]
\item Donner le ratio de couverture après la hausse.
\item Recalculer le spread du modèle après la hausse, la volatilité de l'action étant maintenue à 50~\%.
\item Pourquoi une hausse de l'action et un élargissement du spread peuvent-ils se produire ensemble~?
\item À quel spread de marché l'opération aurait-elle été à l'équilibre sur les deux jours~?
\item Quel instrument aurait couvert la déclassification~?
\end{enumerate}

\medskip\noindent\textbf{Partie IV --- Jugement.}
\begin{enumerate}[resume]
\item Pourquoi de nombreux fonds ont-ils perdu sur des opérations similaires la même semaine~?
\item Comment dimensionner de telles opérations~?
\item Qu'est-ce qui conduirait à clôturer l'opération après les deux jours, et qu'est-ce qui
conduirait à la conserver~?
\item Énoncer le \emph{résultat nommé}~: le P\&L des deux jambes et le total sur les deux jours.
\item En une phrase~: qu'apporte un \omterm{def:rc:structural-credit-models:structural}{modèle structurel} à un desk de crédit~?
\end{enumerate}
\end{problem}

\section{Questions d'entretien}

\begin{interviewq}[$\star$ \iqroles{chercheur, trader}]\label{iq:rc:structural-credit-models:1}
Pourquoi les fonds propres sont-ils une option d'achat sur l'entreprise~? Qu'est-ce que la dette~?
\end{interviewq}

\begin{interviewq}[$\star\star$ \iqroles{chercheur, développeur}]\label{iq:rc:structural-credit-models:2}
Comment estimez-vous la valeur des actifs d'une entreprise et la volatilité de ces actifs~?
\end{interviewq}

\begin{interviewq}[$\star\star$ \iqroles{chercheur}]\label{iq:rc:structural-credit-models:3}
Pourquoi les \omterm{def:rc:structural-credit-models:structural}{modèles structurels} sous-estiment-ils les spreads de crédit à court terme, et comment y remédiez-vous~?
\end{interviewq}

\begin{interviewq}[$\star\star$ \iqroles{risque, banque}]\label{iq:rc:structural-credit-models:4}
Qu'est-ce qu'une \omterm{def:rc:structural-credit-models:dd}{distance au défaut}, et comment la convertit-on en \omterm{def:rc:reduced-form-credit:pd}{probabilité de défaut} en
pratique~?
\end{interviewq}

\begin{interviewq}[$\star\star\star$ \iqroles{trader}]\label{iq:rc:structural-credit-models:5}
Décrivez une opération d'\omterm{def:rc:structural-credit-models:csa}{arbitrage de structure du capital} et son ratio de couverture. Comment peut-elle perdre sur
les deux jambes~?
\end{interviewq}

\begin{interviewq}[$\star\star\star$ \iqroles{chercheur, trader}]\label{iq:rc:structural-credit-models:6}
Comparez les approches à forme réduite et structurelle. Quand utiliseriez-vous chacune~?
\end{interviewq}
